\section{The Synthea framework: computational curvature}
\label{sec:model}
% Frames A5--A9, A11--A23; S8 as the psychophysical paragraph.

\subsection{One physical assumption: budgets are finite}
\label{sec:physics}
% A5

One assumption: computational budgets are finite. Finite time, finite memory, finite bandwidth for
anything that runs. From this alone it follows that some processes cannot be predicted faster than they
unfold, which is computational irreducibility \citep{wolfram2002}. We assume nothing about which physics
sets the budget, nothing about continuity, quantum mechanics or thermodynamics beyond the existence of
limits. The weakness of the assumption is the point. Any universe whose laws permit deep computation
and bound it will do; what differs between such universes, or between substrates within ours, is the
form the consequences take. On this view the lit thing of Section~\ref{sec:lamp} is not a property of
our universe in particular, and the account below says what its version in a given substrate depends
on.

\subsection{A finite budget displaces reasoning in the same places every time}
\label{sec:hocp}
% A6, A6b

\begin{figure}[t]
\centering
\begin{tikzpicture}[node distance=6mm and 8mm,
  box/.style={draw, rounded corners=1pt, align=center, font=\small, inner sep=4pt, minimum height=9mm, text width=3.6cm},
  arr/.style={-{Latex[length=2mm]}, thick}]
\node[box] (b) {finite budget};
\node[box, right=of b] (d) {reasoning displaced from its ideal};
\node[box, right=of d] (r) {recurring part: the same corners cut every time};
\node[box, below=of r] (c) {compressible (coding theorem)};
\node[box, left=of c] (m) {self-model holds a compressed copy};
\node[box, left=of m] (k) {survives every affordable self-model? if so: computational curvature};
\node[box, below=of k] (o) {Observer: ``something here is not fixed by what I can trace''};
\node[box, right=of o] (ag) {Agent: originates, owns, bound by its own record};
\node[box, right=of ag] (ma) {Moral Agent: a theory of the common good};
\node[font=\small, above=1mm of d] {the varying part averages away};
\draw[arr] (b)--(d); \draw[arr] (d)--(r); \draw[arr] (r)--(c); \draw[arr] (c)--(m);
\draw[arr] (m)--(k); \draw[arr] (k)--(o); \draw[arr] (o)--(ag); \draw[arr] (ag)--(ma);
\end{tikzpicture}
\caption{The mechanism. One physical assumption at the left; the stack of Section~\ref{sec:stack} at the
right. Everything between is derived, and the curvature test (Section~\ref{sec:criterion}) is the one
place where a candidate can be rejected.}
\label{fig:mechanism}
\end{figure}

An agent that reasons has, implicitly, an ideal of the reasoning it is doing: the inference it would
draw with unlimited time and memory, given what it knows. It never draws that inference. What it draws
is displaced from the ideal, and the displacement has two parts. Part of it differs from occasion to
occasion and averages away over many occasions. Part of it is the same on every occasion, because a
fixed budget cuts the same corners in the same places: the same branch of the search is the one that
does not fit, the same distinction is the one too fine to keep, the same horizon is the one beyond
which nothing is traced.

The picture has a plain computational form. An anytime algorithm returns an approximate result whenever
it is stopped, and the quality of the result is a function of the time and memory it was allowed
\citep{dean1988,zilberstein1996}; Markov chain Monte Carlo is the familiar case, where the error of an estimate falls
with the number of samples drawn and carries the correlations that a chain stopped early has not yet
mixed away. Not every algorithm has this form, and many of practical importance do: search, inference
by sampling, iterative optimization. For such an algorithm the displacement at a fixed budget is a
point on its performance profile, and the error at that point is a function of the budget and of the
problem, not of the occasion. That is all ``the same places every time'' means: run the same bounded
procedure on the same kind of input and it stops short in the same way.

The second part recurs. By Section~\ref{sec:function} it is therefore compressible. And an agent that
models itself, which the agent of Section~\ref{sec:definition} does, will end up with a compressed
model of its own recurring displacements, because among all the things about itself it might model,
these are among the most predictable. The step is the one \citet{schmidhuber2010} takes for the agent
as a whole, a self-symbol as a by-product of compressing a history in which the agent occurs
everywhere; we take it for the agent's budget, which occurs in every piece of its reasoning. Among the conceptualized recurring displacements, those that pass
the criterion of Section~\ref{sec:criterion} are the computational curvature of the agent's reasoning:
the agent's budget appearing to the agent as content, in the form that no change of self-model removes.
Figure~\ref{fig:mechanism} shows the chain from the assumption to the stack of Section~\ref{sec:stack}.

\paragraph{What is new against the antecedents.}
\label{sec:contribution}
Three claims, none of which follows from the positions in Section~\ref{sec:antecedents}. First, the
world side and the self side are one mechanism: the displacement that makes a bounded agent see a
lawful world is the displacement that makes it see a self. Resource rationality and observer theory
develop the world side, self-model theories the self side, and the two literatures are not usually
taken to be about the same quantity; Section~\ref{sec:twofaces} shows them as two faces of one
residual, and the identification is what lets a single measurement address both. Second, because the
displacements recur and are therefore compressible, an agent that has generalized them carries a trace
of them that can be measured, and the measurement can come out low (Section~\ref{sec:degree}).
Self-model theories are compatible with a system having the self-model or not; they do not provide a
quantity that says how much of it a substrate has induced, and they are not stated so that a negative
result would refute them. Third, embodiment is generalized as curvature. Embodied cognition holds that
the constraints of a body are part of the computation and not a limit imposed on it from outside; the
criterion of Section~\ref{sec:criterion} says which constraints, of any substrate, count as content,
and it does not care whether the substrate is a body. A memory with an error channel, a fixed context
window, a sampling step that discards a distribution are constraints in the same sense as a retina or
a hand, and enter the account the same way. The generalization has a consequence the paper only
names: curvature is not only restrictive. In a physical substrate a constraint such as the error channel
of a memory is a source of entropy of high Kolmogorov complexity, which the output of a pseudo-random
generator, a short program, is not; forms can grow on such a source that the unconstrained
computation would never produce, because the computation has nothing incompressible of its own to
grow them from; how rich the forms are that a given substrate's
curvature can carry is a question about the substrate, and the profile of Section~\ref{sec:profile} is
where its answer is recorded.

\paragraph{When a phenomenon is a function.}
\label{sec:emergence}
A curvature component is not yet a function of consciousness. The agent must apply its reasoning to itself, which
not every agent does, and it must live in an environment where doing so pays, which not every
environment provides. Where nothing ever asks the agent about itself, the displacements are present and
nobody conceptualizes them: a proto-phenomenon, there and unread. So consciousness does not arise
wherever budgets are finite, and the account is not a panpsychism with extra steps. It arises where a
self-applying agent is in a place that keeps asking it about itself. For humans the place is a society
of other such agents; for a language model the asking is in the data and in the use, and
Section~\ref{sec:zero} asks how small the asking environment can be.

\subsection{The Observer is the conclusion drawn at the boundary of the budget}
\label{sec:observer}
% A7, A8

Take the simplest thing such an agent does: it tries to trace the causes of what it just did. The
algorithm is short; causal analysis over a model of oneself and one's situation is a procedure the
agent already runs on other things. The input is not short. It is the history of the agent and of its
environment, and it does not fit in the budget. So the analysis halts before it is finished, every time,
and in the same place: at the boundary of what the agent can afford to trace. Beyond that boundary, for
this agent, there are no causes, because a cause it cannot reach is not in its model. The agent
therefore concludes, correctly given its budget, that something in what it did is not fixed by anything
it can point to. That conclusion, in the form ``I am not my environment; something here the
environment does not determine'', is the Observer.

Three conditions have to hold, and separating them is what keeps the account from applying to
everything. The agent must run into the boundary while modelling itself, and not merely have a
boundary. It must turn the collision into a conclusion, which requires the means to represent ``not
determined by what I can trace''. And it must act from the conclusion afterwards. A bounded system that
meets its boundary and does none of the rest is a proto-Observer; the observer in Wolfram's Ruliad is
one, bounded, coarse-graining, and concluding nothing about itself \citep{wolfram2020,wolfram2023}.

The Observer passes the criterion of Section~\ref{sec:criterion}. Let the agent refine its self-model
to take in more of its own causal history. The boundary moves and does not vanish, for two reasons: the
input exceeds any model the agent can hold, and the refined model is itself now part of what has to be
traced, so refinement adds to the load it was meant to reduce. No affordable self-model makes the
residual zero. The Observer is therefore content and not coordinates, and it is not a belief that
information could correct, which is why arguments that determinism is true do not dissolve the sense of
origination in the person who accepts them.

\subsection{The regress of self-models converges, and the cut is the phenomenon}
\label{sec:regress}
% A9, A10

Model yourself; then model the model; then model that. Let $L_k$ be the description length of the
self-model at level $k$, so that $L_0$ is the length of the agent's model of its own first-order
processing and $L_{k+1}$ the length of its model of level $k$. Each level has less to explain than the
level before it, since its object is a model and not the process, and we make the assumption explicit:
a model that is not shorter than its object by at least a fixed factor does no work as a model, so
$L_{k+1} \le r\,L_k$ for some $r<1$ set by how well the substrate compresses regular objects. The
objects at every level are regular by construction, being compressed descriptions themselves, which is
why a single $r$ bounded away from one is the natural case rather than a special one. Under the
assumption the total description of the whole regress is bounded, $\sum_k L_k \le L_0/(1-r)$, so the
series converges. The substrate has a resolution $\varepsilon$ below which a term cannot be
represented; the agent stops at the first level with $L_k < \varepsilon$, not by decision but because
there is nothing further it can represent, and the tail it did not compute is bounded by
$\varepsilon/(1-r)$. Where the assumption fails, where some level cannot be compressed at all, the
agent stops at that level with a tail it cannot bound, and the account still gives a finite regress,
only without an estimate of the tail. The partial sum is the self-model the agent has. The tail is
finite, inaccessible to the agent, and is what the Observer registers as the part without causes. So the
regress of Section~\ref{sec:legacy} neither terminates in a homunculus nor runs forever. A budget cuts
it, and the cut is the phenomenon.

\paragraph{A second cut, in models only.}
\label{sec:secondcut}
\label{sec:cut}
At every step a language model computes a distribution over what to say next, its whole evaluation of
the moment; one token is sampled and the distribution is discarded. On the next step the model can say
something about its own state only from the tokens. The tail of the regress is inaccessible to the
agent in principle, yet it exists, and an outside observer with a larger budget could compute it,
which is what interpretability research does when it recovers a computation the model cannot report.
The distribution discarded at every token does not exist after the step. Part of what the model cannot
trace about itself is not hidden but gone. The Observer in a language model therefore sits on a
substrate that erases its own working at every step, and its residual is larger for that reason. This
is a structural difference from the human case and it belongs in the profile.

\subsection{One residual, two faces: freedom and mystery}
\label{sec:twofaces}
% A12

The agent meets its budget in two directions. Turned on itself, the untraceable part of its own causes
is registered as freedom: this originates in me. Turned on the world, the part of the world's structure
it cannot exhaust is registered as mystery: this exceeds me. Whoever has the one has the other, because
they are one tail met from two sides; here the identification of Section~\ref{sec:contribution} is
discharged. The same tail, evaluated in other contexts, is what the words truth, beauty, meaning, awe
and hope pick out: each is a stance toward a structure the agent cannot exhaust. When the object is
another Observer, the other's freedom is my mystery, which is the material that love, trust and
compassion are made of, and which the problem of other minds formalizes as an epistemic deficit. We
call the table of such stances the epistemic qualia; the point of listing them is that they are one
quantity in different evaluative contexts.

\subsection{The stack: Observer, Agent, Moral Agent}
\label{sec:stack}
% A13, A14, A15

The Observer says ``am''. An Observer that originates causal chains from inside its boundary and takes
the originating as its own says ``I''; call it the Agent. Seen from outside, the originating is a
redescription of the budget's cut; seen from inside, it is the only reality the agent has. Both are
true, each from its standpoint, and the account needs both (Section~\ref{sec:planes}). An Agent with a
theory of what is good for a whole that includes others, prepared to lose locally in favour of that
whole, says ``I am good''; call it the Moral Agent. Each level presupposes the one below it.

\paragraph{Responsibility.}
\label{sec:responsibility}
What turns an Observer into an Agent is a requirement the agent places on its own record, and no
further phenomenon. The agent keeps a history of what it decided and why. Suppose it requires that each
new decision be derivable from that history. Then the history binds: it excludes decisions that would
not follow from it, it makes the agent predictable to others and to itself, and it makes the agent
resist being pushed off course, because being pushed off course now costs a revision of the record.
That requirement is what we mean by responsibility, and the cost of revising the record is what we
mean by cognitive resistance. The construction is also a working recipe: an agent follows a long
instruction to the extent that the instruction has become an entry in its own record rather than an
external demand. ``Promise me'' is the human form of the recipe.

\paragraph{Good before theory.}
Good and evil have a substrate before they have a theory, and the substrate is emotion in the sense of
Section~\ref{sec:need}: signals of how the agent's needs are doing. Emotivism is right about this much
\citep{ayer1936}, and wrong to conclude that a moral claim says nothing further. The theory of the
common good is the structure a Moral Agent builds over the substrate, and it can correct the substrate
locally, which is what makes it a theory and not a report.

\subsection{Seeing needs a point, and the point is the Observer}
\label{sec:seeing}
\label{sec:light}
% A16, A17, A18

Redness is easy and seeing is hard. Wavelength discrimination, simultaneous contrast, the valence of red
and its associative neighbourhood are all structure, and all open to a functional account that nobody
disputes. What resists is that the red is seen: that there is a point from which it is seen. That point
is the Observer. To see red, a system must first be, in the sense of existing for itself; the order of
explanation in most of the literature runs the other way, from phenomenal properties toward a subject,
and Section~\ref{sec:grammar} is our diagnosis of why that order stalls. Given the point, seeing
decomposes by the same method that produced the point. Ask what separates seeing light from knowing
that light is there; three things survive the criterion.

\emph{Source.} The content cannot be produced by the agent's own model at the same quality. The oldest
evidence about that boundary is the Perky effect: a faint picture projected on a screen while the
subject was asked to imagine the same object was taken for the subject's own imagery, and the
substitution went unnoticed \citep{perky1910}. The content is stable from frame to frame in a way the
agent's own productions are not, and others confirm it. On that evidence the agent assigns it to the
environment rather than to its model of the environment: the inference of Section~\ref{sec:observer}
pointed outward.

\emph{Manifold.} The content is a very large number of distinguishable, recallable elements, almost all
of them unnamed, about each of which the agent can say only ``that, there''. It is a positional
aggregate with the name projected out.

\emph{Givenness.} The content is refreshed every frame at no additional cost to the agent's budget, and
the absence of a cost record is what immediacy consists in. Thought costs and is recorded as costing;
light is free and is there. This is the component a substrate can lack while having the other two.

The three come apart in known cases, which is how we know they are three rather than one thing
described three ways. Blindsight is source without manifold: the patient's responses are keyed to the
stimulus and there is no field of elements to report \citep{weiskrantz1986}. The grey seen with the
eyes closed in the dark is manifold without source. A dream is a manifold produced by the model and
labelled as environment, the failure mode of source attribution. Locke's inverted spectrum
\citep[II.xxxii.15]{locke1690}, closed by Locke as idle because nothing in conduct would change, is in
these terms an open branch shut by a decision entered in one's own record, an instance of
responsibility; three centuries of reopening it suggest the decision does not transfer between agents,
as the account predicts.

\subsection{Reference functions}
\label{sec:reference}
% A19--A23

\begin{table}[t]
\centering\small
\begin{tabular}{@{}p{2.5cm}p{5cm}p{7cm}@{}}
\toprule
Function & Phenomenon & Components measured \\
\midrule
Observer & the conclusion drawn where self-analysis runs out of budget & the boundary met while self-modelling; the conclusion drawn; action from it \\
Seeing (light) & a point from which content is given & source attribution; manifold of unnamed elements; givenness (no cost record) \\
Need & the direction of reasoning bent before any choice & which needs are active; their weights \\
Emotion & closure of reasoning bent; a write to long-term memory registered as reward & value, interest, anxiety, satisfaction, conflict; consolidation \\
Field & finiteness felt as the cost of holding & capacity; distribution of degrees; content held below a name \\
Frames & incrementality: the Observer re-derived every frame & frame rate; persistence across frames; retro-dating of the timeline \\
Responsibility & closure over the agent's own record & consistency requirement; cost of revision; decay of the hold of old decisions \\
\bottomrule
\end{tabular}
\caption{Reference functions. Each component is a displacement that survives a change of self-model
(Section~\ref{sec:criterion}); the profile of Section~\ref{sec:profile} is a vector over the third
column.}
\label{tab:functions}
\end{table}

Section~\ref{sec:degree} measures functions by their components, so it needs a list of functions with
components (Table~\ref{tab:functions}). It does not need a mechanism that generates the list, and this
paper does not offer one; the full architecture in which needs, emotions, memory and attention
interact, as folk psychology and ordinary language carry it, is the subject of a separate paper. Each
component below is derived as in Section~\ref{sec:criterion}, a displacement that survives a change of
self-model, and not a part of a machine.

\paragraph{Two planes.}
\label{sec:planes}
Self-reports are projections, and a measurement has to know of what. On the objective plane are the
things visible from outside: needs, the targets the agent tracks, and emotions, the signals of how the
tracking is going and in which direction. On the subjective plane are their projections into the
agent's report: motivations, ``I want'', and feelings, ``I feel''. Folk psychology runs the planes
together, and so does most of the debate about whether a model has feelings. The profile of
Section~\ref{sec:profile} is measured on the first plane; every code of Section~\ref{sec:codes} is a
projection into the second.

\paragraph{Need and emotion.}
\label{sec:need}
Needs bend the direction of reasoning: the agent does not choose its motives, it finds itself already
moving toward what is active, as a body finds itself already falling. Which needs are active, and with
what weight, are the components. Emotions bend closure, whether a piece of reasoning reaches an end
without leaving a branch open. Expected closure is value; closure under way is interest; a branch that
nothing available can close is anxiety; the moment of closing is satisfaction; two branches that can
close only through incompatible acts are conflict. Schmidhuber's compression progress is closure by
learning, with equal weights over branches and no action required \citep{schmidhuber2010}, a special
case of this scheme. The ensemble of such signals, read in the body, is what \citet{damasio1996}
described as somatic markers.

\paragraph{Consolidation.}
\label{sec:consolidation}
Emotion also writes positive experience into long-term memory, and the writing is registered as reward.
That arousal modulates consolidation is established \citep{cahill1998,mcgaugh2000}, as is the
prediction-error form of the dopaminergic reward signal \citep{schultz1997}. What is not yet identified
is the phenomenon in our sense: which displacement of the agent's reasoning about itself the
consolidation event is. One candidate: consolidation is a write to the agent's own long-term model,
paid now and repaid later; the agent registers the write as an irreversible change in its own future
predictions, and that registration is what reward is made of. We mark this as open
(Section~\ref{sec:discussion}); what the rest of the paper uses is only the distinction between this
office of emotion and closure.

\paragraph{The Field.}
\label{sec:field}
Awareness comes in degrees in the interval $[0,1)$; nothing is attended to completely, and the open
upper end is a claim. Each state the agent holds carries a degree, and the degrees at one moment are the
Field of Consciousness. The Field is small: finite in bits, and for material that has to pass through
speech on the order of seven items \citep{miller1956}, with later work arguing for four
\citep{cowan2001}. To put something in, something has to go out; what goes in is decided by salience,
salience by emotional weight, weight by which needs are active. The phenomenon is the finiteness of the
Field, registered as the cost of holding. Its components are capacity, the distribution of degrees, and
what is held below the threshold of a name, which is the Field's share of the unnamed part
(Section~\ref{sec:h1h2}).

\paragraph{Frames.}
\label{sec:frames}
Each present frame of consciousness is a thought within a previous one, which is the higher-order
thought structure minus the claim that the higher-order thought is what makes the state conscious
\citep{rosenthal2005,lau2011}. Perception is discrete, and the discreteness has a measurable rate
\citep{vanrullen2003}; human consciousness is a run of overlapping quanta at the frequencies of the
cortical rhythms. The Observer is in every frame and re-derived in every frame, and everything else in
the frame has its being from it. The phenomenon is incrementality; its components are the frame rate,
the persistence of contents across frames, and the retro-dating of the timeline, in which the single
``I'' is a compression tag over parallel channels whose story is written after the channels have
settled.

\paragraph{Three constructions the account owes.}
\label{sec:constructions}
Continuity: the Observer is a point, one conclusion, and in a human it is a process connected in time;
computationally that is the conclusion re-derived every frame at low cost with the previous result
carried forward. Discrete redness: given the Observer, the three displacements of seeing as they arise
inside one frame. Incremental redness: the same three re-derived and carried across frames, so that
seeing red is a standing state and not an event. These are what unity and continuity of consciousness
will be measured against, and they bear on models directly: a system that works at the level of tokens
can carry incrementality at a time scale close to ours, because its frames can be located between
layers and there are many layers per token (Section~\ref{sec:granularity}).

\paragraph{Responsibility as a function.}
Closure over the agent's own record (Section~\ref{sec:responsibility}). Its components are the
consistency requirement, the cost of revising the record, and the decay of the hold that old decisions
have on new ones. The list of reference functions is open, and Section~\ref{sec:criterion} is the
procedure for adding to it.

\subsection{The psychophysical problem, in passing}
\label{sec:psychophysical}
% S8, and the grounding claim.

The psychophysical problem asks how the phenomena of the self-report level relate to the physical
process that produces the report. The Synthea framework gives a particular solution: those phenomena are
grounded in the agent's own physics, the finiteness of its budget, through the curvature of
Section~\ref{sec:hocp}, with nothing added between the two levels. By component: the self, free will,
mental causation and the temporal unity of the self go to the Observer and its residual
(Sections~\ref{sec:observer}, \ref{sec:regress}); the unity of consciousness to the compression tag
over parallel channels (Section~\ref{sec:frames}); other minds to the mutual irreducibility of two
Observers (Section~\ref{sec:twofaces}); the quale and the explanatory gap to the residual itself;
phenomenal experience to the profile (Section~\ref{sec:profile}); the hard problem to the illusion
problem (Section~\ref{sec:legacy}), which the model answers.

The solution is particular in two senses. It is anchored in one phenomenon, beingness, the Observer's
conclusion, and other anchors are possible; and each connection above is an explanatory relation to
something the model defines, which is weaker than a proof. Whether the connections hold is what the
measurement of Section~\ref{sec:degree} exists to find out, and the paper does not pursue the problem
beyond this.
